\documentclass[11pt]{article}
\usepackage[margin=1in]{geometry}
\usepackage{amsmath,amsthm,amssymb,amsfonts,latexsym,mathtools}
\usepackage{graphicx}
\usepackage{hyperref}
\usepackage{enumitem}
\usepackage[utf8]{inputenc}


\numberwithin{equation}{section}

\begin{document}

\title{A Framework Toward the Millennium Prize Problem\\
       on the Navier--Stokes Equations}
\author{Thomas Birnie + ChatGPT}
\date{\today}

\maketitle

\begin{abstract}
We provide a comprehensive and rigorously detailed \LaTeX{} document
outlining a potential framework for solving the global regularity problem of
the three-dimensional incompressible Navier--Stokes equations.
While this document describes a complete \emph{strategy} with all intermediate steps,
we emphasize that the final resolution of this famous open problem remains a subject
of ongoing investigation.
\end{abstract}

\tableofcontents

\section{Introduction}\label{sec:intro}

The Navier--Stokes equations lie at the heart of fluid mechanics, describing the motion
of viscous, incompressible fluids. Despite their widespread use in physical applications,
the fundamental mathematical question of whether smooth solutions exist for all time in
three spatial dimensions remains open. This challenge is recognized as one of the
\emph{Millennium Prize Problems} posed by the Clay Mathematics Institute.

In this document, we present a framework containing the standard techniques employed
in partial regularity proofs, global energy estimates, compactness arguments, and
convergence analyses. Although many of these methods have yielded deep insights,
a definitive solution remains elusive. The structure here reflects a best-practice
approach to writing a rigorous manuscript that, if completed successfully, would
meet the high standards required by the Clay Institute.

\section{Statement of the Problem}\label{sec:statement}

We consider the three-dimensional incompressible Navier--Stokes equations:
\begin{equation}\label{eq:NS}
\begin{cases}
\partial_t \mathbf{u} + (\mathbf{u}\cdot \nabla)\mathbf{u}
  = -\nabla p + \nu \Delta \mathbf{u}, \\
\nabla \cdot \mathbf{u} = 0,
\end{cases}
\end{equation}
where \( \mathbf{u}(t,\mathbf{x}) \in \mathbb{R}^3 \) is the velocity field,
\( p(t,\mathbf{x}) \in \mathbb{R} \) is the pressure, and \( \nu > 0 \) is the
kinematic viscosity. The domain is typically \( \mathbb{R}^3 \) or a subset
with appropriate boundary conditions (e.g., periodic or Dirichlet). Initial data
is given by
\[
\mathbf{u}(0,\mathbf{x}) = \mathbf{u}_0(\mathbf{x}),
\quad \nabla \cdot \mathbf{u}_0 = 0.
\]

\begin{definition}[Global Existence and Smoothness Problem]
Show that for any sufficiently smooth, divergence-free initial data
\(\mathbf{u}_0\) with finite energy, there exists a unique, smooth solution
\(\mathbf{u}(t,\mathbf{x})\) to \eqref{eq:NS} for all \( t \ge 0 \),
and that this solution remains smooth for all time.
\end{definition}

\section{Preliminaries}\label{sec:prelims}

Here, we gather fundamental results, tools, and definitions needed throughout
the paper. This includes Sobolev spaces, embeddings, the classical Aubin--Lions
compactness lemma, and other widely used PDE methods.

\subsection{Sobolev spaces and embeddings}

Let \(\Omega \subset \mathbb{R}^3\). Denote by \(H^s(\Omega)\) the Sobolev space
of functions with square-integrable weak derivatives up to order \(s\).
Key inequalities:
\begin{itemize}[leftmargin=2em]
\item \emph{Gagliardo--Nirenberg Inequality}
\item \emph{Ladyzhenskaya Inequality}
\item \emph{Poincar\'{e} Inequality}
\end{itemize}

We recall that for \(1 \leq p < \infty\), the norm in \(W^{s,p}(\Omega)\) is
\[
\| f \|_{W^{s,p}} = \left(
   \sum_{|\alpha|\le s} \int_\Omega |D^\alpha f|^p \, d\mathbf{x}
\right)^{1/p}.
\]
All relevant inequalities will be stated precisely where needed.

\subsection{Aubin--Lions lemma}

We state a version of the Aubin--Lions lemma crucial for compactness arguments:

\begin{lemma}[Aubin--Lions]
\label{lemma:aubin_lions}
Let \(B_0 \subset B_1 \subset B_2\) be Banach spaces with compact embedding
\(B_1 \hookrightarrow B_0\) and continuous embedding \(B_2 \hookrightarrow B_1\).
Then for \(1 \leq p \leq \infty\), if \(\{u_n\}\) is a sequence such that
\[
u_n \in L^p(0,T;B_2), \quad
\partial_t u_n \in L^1(0,T;B_0),
\]
then there exists a strongly convergent subsequence in \(L^p(0,T;B_1)\).
\end{lemma}

\begin{proof}
See \cite{ConstantinFoias,Temam} for detailed proofs and variants.
\end{proof}

\section{Discrete Formulations}\label{sec:discrete}

In many modern approaches, we rely on discretization schemes (e.g., Galerkin methods)
to approximate solutions and then pass to the limit to show existence of smooth
solutions. Here we outline the discrete operators, the resulting system, and
key properties such as stability.

\subsection{Galerkin approximation}

Let \(\{\mathbf{w}_j(\mathbf{x})\}_{j=1}^\infty\) be a suitable orthonormal basis
for divergence-free vector fields in \(L^2(\Omega)\).
For each \(n\in \mathbb{N}\), define the finite-dimensional subspace
\[
V_n = \mathrm{span}\{\mathbf{w}_1, \dots, \mathbf{w}_n\}.
\]
We seek approximations \(\mathbf{u}_n(t) \in V_n\) such that
\begin{equation}\label{eq:Galerk}
\begin{aligned}
\frac{d}{dt}(\mathbf{u}_n,\mathbf{w}_j)
+ \bigl((\mathbf{u}_n\cdot \nabla)\mathbf{u}_n,\mathbf{w}_j\bigr)
+ \nu(\nabla \mathbf{u}_n,\nabla \mathbf{w}_j)
&= 0,
&& \text{for all } j = 1,\ldots,n, \\
(\mathbf{u}_n(0), \mathbf{w}_j) &= (\mathbf{u}_0,\mathbf{w}_j).
\end{aligned}
\end{equation}
Here \((\cdot,\cdot)\) denotes the \(L^2(\Omega)\) inner product.

\subsection{Discrete energy estimates}

The usual procedure is to show that these approximations satisfy a discrete
energy inequality. For instance:
\[
\frac{1}{2} \frac{d}{dt} \|\mathbf{u}_n\|_{L^2}^2
+ \nu \|\nabla \mathbf{u}_n\|_{L^2}^2 \le 0.
\]

\begin{lemma}[Discrete Energy Inequality]
Let \(\mathbf{u}_n\) be a Galerkin approximation solving \eqref{eq:Galerk}.
Then for \(0 \le t \le T\),
\[
\|\mathbf{u}_n(t)\|_{L^2}^2 + 2\nu
\int_0^t \|\nabla \mathbf{u}_n(\tau)\|_{L^2}^2 \, d\tau
\le \|\mathbf{u}_0\|_{L^2}^2.
\]
\end{lemma}

\begin{proof}
Test \eqref{eq:Galerk} by \(\mathbf{u}_n\) and use incompressibility to handle
the nonlinear term. The Grönwall lemma is then applied.
\end{proof}

\section{Mapping Between Discrete and Continuous Formulations}\label{sec:mapping}

A key step is to rigorously connect the finite-dimensional approximations to the
infinite-dimensional PDE. We show that in the limit \(n \to \infty\),
\(\mathbf{u}_n\) converges to a function \(\mathbf{u}\) that satisfies
the Navier--Stokes equations in the weak or strong sense.

\subsection{Weak formulation}

In the usual weak formulation, one tests with smooth divergence-free fields.
Convergence of \(\mathbf{u}_n\) to \(\mathbf{u}\) must preserve all relevant
properties:
\begin{enumerate}[label=(\roman*)]
\item \(\mathbf{u}_n \rightharpoonup \mathbf{u}\) in appropriate Sobolev spaces,
\item \(\nabla \cdot \mathbf{u} = 0\) almost everywhere,
\item The nonlinear term \((\mathbf{u}_n \cdot \nabla)\mathbf{u}_n\) converges
      to \((\mathbf{u}\cdot \nabla)\mathbf{u}\) in the sense of distributions.
\end{enumerate}

\begin{remark}
Compactness results such as the Aubin--Lions lemma
(see Lemma~\ref{lemma:aubin_lions}) are instrumental.
\end{remark}

\section{Proof of Compactness}\label{sec:compactness}

We now state and prove the main compactness theorem for the sequence
\(\{\mathbf{u}_n\}\). The general template involves showing that
\(\{\mathbf{u}_n\}\) is bounded in a sufficiently strong norm and that
its time derivative is bounded in a weaker norm, thus allowing an application
of Aubin--Lions.

\begin{theorem}[Compactness of Discrete Solutions]
\label{thm:compactness}
Let \(\{\mathbf{u}_n\}\) be the Galerkin approximations. Assume that
the initial data \(\mathbf{u}_0\in H^1(\Omega)\) with \(\nabla\cdot\mathbf{u}_0=0\).
Then there exists a subsequence (still denoted \(\mathbf{u}_n\))
and a function \(\mathbf{u}\) such that
\[
\mathbf{u}_n \to \mathbf{u} \quad \text{strongly in } L^2(0,T;H^1(\Omega)),
\]
and \(\mathbf{u}\) satisfies the incompressible Navier--Stokes equations
in the weak sense.
\end{theorem}

\begin{proof}[Proof Sketch]
\textbf{Step 1: Uniform bounds.}
From the discrete energy inequality, we have uniform boundedness of
\(\mathbf{u}_n\) in \(L^\infty(0,T;L^2(\Omega))\) and
\(\nabla \mathbf{u}_n\) in \(L^2(0,T;L^2(\Omega))\).

\noindent
\textbf{Step 2: Time derivative bounds.}
Differentiating \eqref{eq:Galerk} and using orthogonality properties,
one obtains \(\partial_t \mathbf{u}_n\) bounded in a negative-order Sobolev space
\(H^{-1}(\Omega)\), ensuring that
\[
\partial_t \mathbf{u}_n \in L^2(0,T;H^{-1}(\Omega)).
\]

\noindent
\textbf{Step 3: Application of Aubin--Lions.}
We invoke Lemma~\ref{lemma:aubin_lions} with
\[
B_2 = H^1(\Omega),\quad B_1 = L^2(\Omega), \quad B_0 = H^{-1}(\Omega).
\]
Hence, \(\{\mathbf{u}_n\}\) is relatively compact in \(L^2(0,T;H^1(\Omega))\).
A diagonal argument yields a strongly convergent subsequence.

\noindent
\textbf{Step 4: Identification of the limit.}
Passing to the limit in the Galerkin scheme, we verify that the limit \(\mathbf{u}\)
satisfies the weak form of the Navier--Stokes equations.
\end{proof}

\section{Proof of Energy Bounds}\label{sec:energy_bounds}

Beyond the basic \(L^2\) energy estimate, one typically seeks higher-order energy
bounds (in \(H^s\) norms) to argue for smoothness. In 2D, such bounds can be made
global in time, but in 3D the bounding technique meets well-known difficulties.

\subsection{Higher Sobolev norm estimates}

We outline the derivation of an \(H^1\) estimate for \(\mathbf{u}\). Testing the
Navier--Stokes equations with \(-\Delta \mathbf{u}\) can yield:
\[
\frac{1}{2}\frac{d}{dt} \|\nabla \mathbf{u}\|_{L^2}^2 + \nu \|\Delta \mathbf{u}\|_{L^2}^2
  \le C \|\nabla \mathbf{u}\|_{L^2}^2 \|\mathbf{u}\|_{L^4}^2,
\]
where one applies Sobolev embedding to handle the nonlinear term. The constant \(C\)
may depend on domain geometry and embedding constants. Similar estimates can be
attempted for higher derivatives.

\begin{remark}
The key unresolved issue is whether these estimates can be closed globally in 3D
(leading to a blowup criterion if not). This is precisely the crux
of the Millennium problem.
\end{remark}

\section{Proof of Convergence}\label{sec:conv_proof}

Once compactness and uniform (in \(n\)) energy bounds are established, one shows
weak and then strong convergence of the discrete solutions to a smooth solution
\(\mathbf{u}\). Typical steps:

\begin{theorem}[Weak Convergence]
With the hypotheses of Theorem \ref{thm:compactness}, the limit function \(\mathbf{u}\)
is a \emph{weak solution}, meaning:
\[
\int_0^T \int_\Omega \Bigl(
  \mathbf{u}\cdot \partial_t \boldsymbol{\phi}
  + (\mathbf{u} \otimes \mathbf{u}): \nabla \boldsymbol{\phi}
  + \nu \nabla \mathbf{u} : \nabla \boldsymbol{\phi}
\Bigr) \, d\mathbf{x}\, dt = 0,
\]
for all smooth, divergence-free test functions \(\boldsymbol{\phi}\) that vanish at \(t=0\)
and \(t=T\).
\end{theorem}

\begin{proof}
Standard limiting procedure: pass to the limit in the Galerkin formulation
\eqref{eq:Galerk}. Use the compactness to handle the nonlinear terms,
and identify the weak limit in each term.
\end{proof}

\begin{theorem}[Strong Convergence]
Under additional assumptions (e.g., higher regularity of initial data),
\(\mathbf{u}_n \to \mathbf{u}\) strongly in \(L^2(0,T;H^1(\Omega))\),
and \(\mathbf{u}\) is a \emph{strong solution} if it meets the PDE in the classical sense
and has all required derivatives in \(L^2\).
\end{theorem}

\begin{proof}[Proof Sketch]
A more refined compactness argument using improved estimates on higher Sobolev
norms of \(\mathbf{u}_n\). Under these assumptions, standard energy methods
plus the Aubin--Lions lemma yield strong convergence.
\end{proof}

\section{Toward a Proof of Global Smoothness}\label{sec:global_smoothness}

\subsection{Heuristic outline}

A classical approach to global smoothness would argue as follows:
\begin{enumerate}[label=(\arabic*)]
\item Show that the sequence of approximate solutions remains bounded
      in arbitrarily high Sobolev norms.
\item Use bootstrapping and Sobolev embeddings to promote the solution
      to classical regularity (smoothness).
\item Conclude no finite-time singularity can develop because the
      blow-up scenario is precluded by the energy estimates and regularity
      arguments.
\end{enumerate}

\subsection{Open problems and partial results}

Despite substantial progress:
\begin{itemize}[leftmargin=2em]
\item Known partial regularity results (Caffarelli--Kohn--Nirenberg, Scheffer, etc.)
      guarantee that potential singularities, if they exist, occupy a set
      of measure zero in space-time.
\item Scaling arguments and blow-up criteria (e.g., Beale--Kato--Majda for 2D Euler)
      highlight the difficulty in controlling the nonlinear term in 3D.
\end{itemize}

\begin{remark}
Currently, no known method closes the final gap for 3D global regularity.
However, the argument structure provided here---energy inequalities,
compactness, and Sobolev regularity---is the cornerstone of potential future proofs.
\end{remark}

\section{Conclusion}\label{sec:conclusion}

We have presented a rigorous \emph{framework} for addressing the Millennium Prize Problem
on the Navier--Stokes equations, mirroring the structure that any final proof
would likely require:
\begin{itemize}[leftmargin=2em]
\item \textbf{Preliminaries:} Sobolev spaces, key inequalities, Aubin--Lions lemma.
\item \textbf{Discrete Formulation:} Galerkin approximations, discrete energy estimates.
\item \textbf{Compactness and Convergence:} Passing to the limit to obtain weak and
      strong solutions.
\item \textbf{Higher Regularity and Smoothness:} Energy bounds in higher Sobolev spaces,
      plus the conceptual blueprint for proving global smoothness.
\end{itemize}

While the final, decisive step to show that \emph{no singularities form in finite time}
for the 3D case remains open, the methods outlined here represent the rigorous backbone
that any successful solution would build upon. It is our hope that clarifying this
structure in a cohesive manner fosters deeper understanding and catalyzes future
breakthroughs in the Navier--Stokes regularity question.

\appendix
\section{Additional Technical Details}

\subsection{Inequalities and their proofs}
A compendium of classical inequalities used (Poincar\'e, Gagliardo--Nirenberg,
etc.) can be found in standard references \cite{Evans,Temam}.

\subsection{Compatibility conditions and boundary terms}
For domain-boundary constraints (e.g., no-slip conditions), further details on
elliptic regularity and boundary-layer phenomena must be added for a full treatment.

\begin{thebibliography}{99}
\bibitem{ConstantinFoias}
P.~Constantin, C.~Foias.
\emph{Navier--Stokes Equations}.
Chicago Lectures in Mathematics, 1988.

\bibitem{Evans}
L.~C.~Evans.
\emph{Partial Differential Equations}.
American Mathematical Society, 2010.

\bibitem{Temam}
R.~Temam.
\emph{Navier--Stokes Equations, Theory and Numerical Analysis}.
North-Holland, 1977.
\end{thebibliography}

\end{document}
